% ============================================================
%  macros.tex —— 全书公共宏定义（由 Lead 维护，撰稿人请勿修改）
%  使用前已加载：amsmath, amssymb, bm, booktabs, tikz, tcolorbox,
%  listings, enumitem, pgfplots, hyperref, adforn
% ============================================================

% ---------- 数学快捷命令 ----------
\newcommand{\R}{\mathbb{R}}
\newcommand{\SO}{\mathrm{SO}}
\newcommand{\SE}{\mathrm{SE}}
\newcommand{\so}{\mathfrak{so}}
\newcommand{\se}{\mathfrak{se}}
\newcommand{\norm}[1]{\left\lVert #1 \right\rVert}
\newcommand{\abs}[1]{\left\lvert #1 \right\rvert}
\newcommand{\diff}{\mathrm{d}}
\newcommand{\E}{\mathbb{E}}
\DeclareMathOperator{\tr}{tr}
\DeclareMathOperator{\diag}{diag}
\DeclareMathOperator{\Cov}{Cov}
\DeclareMathOperator{\Var}{Var}
\DeclareMathOperator*{\argmin}{arg\,min}
\DeclareMathOperator*{\argmax}{arg\,max}
\newcommand{\degree}{^\circ}
\newcommand{\eg}{e.g.}
\newcommand{\ie}{i.e.}

% ---------- 本章导读（仿 ElegantBook introduction 样式）----------
\newtcolorbox{introbox}{
  enhanced,
  breakable,
  colback=structurecolor!7,
  colframe=structurecolor,
  boxrule=0pt,
  toprule=0.8pt,
  bottomrule=0.8pt,
  sharp corners,
  colbacktitle=structurecolor,
  coltitle=white,
  fonttitle=\bfseries,
  title={本章导读},
  attach boxed title to top center={yshift=-3mm,yshifttext=-1mm},
  boxed title style={sharp corners, boxrule=0pt},
  top=9pt, bottom=7pt, left=3mm, right=3mm,
  before skip=10pt, after skip=10pt
}
\newcommand{\chapterintro}[1]{\begin{introbox}#1\end{introbox}}

% ---------- 要点框 ----------
\newtcolorbox{keybox}[1]{
  enhanced,
  breakable,
  colback=structurecolor!4,
  colframe=structurecolor!55,
  fonttitle=\bfseries\small,
  colbacktitle=structurecolor!12,
  coltitle=structurecolor!80!black,
  title={#1},
  boxrule=0.5pt,
  arc=1.5mm,
  left=2mm, right=2mm, top=1.2mm, bottom=1.2mm,
  before skip=6pt, after skip=6pt,
  drop shadow={opacity=0.10}
}

% ---------- 伪代码环境 ----------
\newtcolorbox{pseudobox}{
  enhanced,
  breakable,
  colback=white,
  colframe=black!25,
  boxrule=0.5pt,
  arc=1mm,
  left=2mm, right=2mm, top=1.5mm, bottom=1.5mm,
  drop shadow={opacity=0.08}
}
\newenvironment{pseudo}[1]{%
  \begin{pseudobox}%
  \noindent\textbf{\color{structurecolor}算法：}\textsf{#1}\par\vspace{1mm}%
  \begin{enumerate}[label=\textbf{\arabic*.}, leftmargin=*, topsep=1pt, itemsep=0.5pt, parsep=0pt]%
}{%
  \end{enumerate}%
  \end{pseudobox}%
}

% ---------- 章末练习（花饰标题，不写入目录）----------
\renewenvironment{problemset}[1][\xchaptertitle~\exercisename]{%
  \vspace*{10pt}%
  \begin{center}%
    \markright{#1}%
    \textcolor{structurecolor}{\Large\bfseries
      \adftripleflourishleft~#1~\adftripleflourishright}%
  \end{center}%
  \begin{enumerate}%
}{%
  \end{enumerate}%
}

% ---------- 本章小结（花饰标题）----------
\newenvironment{summary}{%
  \par\vspace*{12pt}%
  \begin{center}%
    \textcolor{structurecolor}{\Large\bfseries
      \adftripleflourishleft~本章小结~\adftripleflourishright}%
  \end{center}%
  \vspace{-2pt}%
  \noindent\ignorespaces
}{%
  \par\vspace*{4pt}%
}

% ---------- 延伸阅读（花饰标题）----------
\newenvironment{reading}{%
  \par\vspace*{10pt}%
  \begin{center}%
    \textcolor{structurecolor}{\large\bfseries
      \adftripleflourishleft~延伸阅读~\adftripleflourishright}%
  \end{center}%
  \vspace{-2pt}%
  \noindent\ignorespaces
}{%
  \par\vspace*{2pt}%
}

% ---------- 常用 TikZ 样式（可直接在图中使用）----------
\tikzset{
  block/.style   = {rectangle, rounded corners=2pt, draw=structurecolor,
                    thick, fill=structurecolor!8, minimum height=8mm,
                    minimum width=18mm, align=center, font=\small},
  block2/.style  = {rectangle, rounded corners=2pt, draw=structurecolor!70,
                    thick, fill=structurecolor!18, minimum height=8mm,
                    minimum width=18mm, align=center, font=\small},
  arrow/.style   = {-{Stealth[length=2mm]}, thick, draw=black!70},
  darrow/.style  = {{Stealth[length=2mm]}-{Stealth[length=2mm]}, thick,
                    draw=black!70},
  label/.style   = {font=\footnotesize, inner sep=1pt},
}
